% JNE-specific replacement of original BSPC editorial comparison.
\subsection{Relation to established neural-engineering validation studies}
\label{sec:p1-jne-peer}
Recent accepted research in Journal of Neural Engineering shows
that a neural decoding methods paper need not report a new SOTA
network, but does need a generalisable, experimentally resolved
measurement or methodological question.
Eder et al.\ \cite{eder2024jnebenchmark} compared Riemannian
decoding against EEGNet and other convolutional models
across motor imagery, P300 and SSVEP paradigms with multiple
validation protocols and reproducible open-source pipelines.
Seibel et al.\ \cite{seibel2026jnedata}
measured how CNN decoding performance responds to subject
count and recorded EEG duration across three tasks.
Their conclusions are operational: they directly inform
model family or data acquisition decisions.
Our scope is narrower---binary motor-imagery
tangent-space linear discriminants and the
selection \emph{criterion}, not a deep learning
architecture leaderboard nor a data-acquisition
learning curve.

\begin{table}[tb]
\centering\small
\caption{Comparison with representative published
JNE methods studies, preserving non-comparable
endpoints. This is a design/evidence comparison,
not an accuracy leaderboard.}
\label{tab:p1-jne-peer}
\setlength{\tabcolsep}{4pt}
\begin{tabular}{p{.24\linewidth}p{.31\linewidth}p{.33\linewidth}}
\toprule
Study & Experimental scope & Actionable conclusion\\
\midrule
Eder et al.\ \cite{eder2024jnebenchmark}
& MI/P300/SSVEP; CNN and Riemannian; multiple split budgets
& Comparative model-family tradeoffs\\
Seibel et al.\ \cite{seibel2026jnedata}
& Three EEG tasks and CNN structures; varied training sample supply
& Estimate marginal data-acquisition value\\
This study
& Four original frozen AUC experiment strata; three
  exposed source-only EEG score-anatomy collections
& Separate genuinely predictive from class-centred
  source scoring and disclose invariant penalties;
  no new optimal transfer selector established\\
\bottomrule
\end{tabular}
\end{table}

We aim to supply an equally \emph{verifiable} conclusion
within our constrained decoder family: a historical
validation covariance score can be misidentified as
source-trained class-conditional predictive likelihood,
and even a correctly computed Gaussian score can be
dominated by source validation shifts invisible to
a fixed within-person rank metric.
The scientific value lies in the distinction
between these estimands and the downstream
interpretation of historical source choices.
Unlike the two cited JNE examples, we have
\emph{not} demonstrated that a replacement
scoring policy gives consistently better
future-participant decoding. That limitation
is an explicit boundary of the manuscript.

\subsection{Outcome consequences after correcting the comparator}
\label{sec:p1-correct-comparator-decision}
The original pre-data comparison of source-macro
AUC selection against the historical class-centred
CovScore is an authentic outcome comparison but
not a test against properly predictive Gaussian
model selection. The latter evaluation was
specified after earlier outcomes had been read.
To prevent an invalid upgrade of the original
test, we report the corrected post-outcome
matched contrasts separately:
\begin{table}[tb]
\centering\small
\caption{Proper predictive source-selection
comparisons on already exposed EEG results.
These are \emph{descriptive exploratory}
target outcomes, not two new confirmatory wins.
Person-macro ROC-AUC values must not be
compared directly with accuracy percentages
from different EEG paradigms.}
\label{tab:p1-correct-predictive}
\begin{tabular}{lrrr}
\toprule
Cohort & Macro source AUC & Proper NLL AUC &
Macro minus proper\\
\midrule
Shin2017A ($n=29$) &0.785900&0.786207&$-0.000307$\\
BNCI2014\_002 ($n=14$) &0.778705&0.771886&$+0.006819$\\
\bottomrule
\end{tabular}
\end{table}
For the 14-person BNCI2014\_002 comparison,
the exploratory interval for the final contrast
was $[-0.00595,0.02027]$, containing zero.
The Shin2017A corrected comparator performed
slightly \emph{better} than source-macro
selection. Both observations are counterevidence
to any strong claim that one selection objective
is universally superior. The source-selected
shrinkage choices also differ between acquisitions,
and the large fixed-model density-score
penalties cannot settle the sign of an
unseen-person AUC difference.
The practical recommendation is therefore a
\emph{measurement and reporting protocol},
not an unvalidated replacement algorithm:
\begin{enumerate}
\item Label each covariance-validation
metric by the origin of every class mean
(train-source versus validation-subject).
\item State whether held source labels enter
the evaluation score, and distinguish such
source labels from prohibited held-target
information.
\item Report paired held-person rank metrics
and negative effects independently of
the density-score change and model-family
choice.
\item Distinguish pre-outcome source
decisions from any post-outcome diagnostic,
even if both use public EEG data.
\item Keep constant the sensor configuration,
training source participants, target information
budget and original frozen parameter grid
in any algorithmic superiority comparison.
\end{enumerate}
The exact three-factor decomposition can serve
as a diagnostic explaining density-score
variance in this particular linear family.
It is not proven to forecast whom to calibrate,
whether to abstain, or which model to deploy
on an unseen individual.
